\section{Self-adjointization, coupling, and controls}
\label{sec:controls}

The most natural self-adjoint block erases rather than reveals Haar phase.
Set
\[
 H=\begin{pmatrix}0&W\\W^*&0\end{pmatrix}.
\]
Then
\[
 H^2=I.                                                    \tag{5.1}
\]
Its spectrum is contained in $\{-1,1\}$, independent of the circle-valued
phase of $u$.  The weighted block
\[
 H_q=\begin{pmatrix}0&qW\\\overline qW^*&0\end{pmatrix}
\]
satisfies $H_q^2=|q|^2I$ and introduces $\overline q$, so it is not the frozen
holomorphic $s$-determinant.

Recurrent symbolic coupling makes the phase trace-visible only by creating
the wrong closed words.  Consider two vertices joined by directed edges with
formal scalar weights $x,y$ and fiber labels $u,u^{-1}$:
\[
 \cL=\begin{pmatrix}0&xu\\yu^{-1}&0\end{pmatrix}.         \tag{5.2}
\]
The balanced words $uu^{-1}=u^{-1}u=e$ give
\[
 \Phi_c^{(2)}(\cL^2)=2(1+c)xy.                             \tag{5.3}
\]
This is a mixed two-cycle.  It is neither a pure $p^2$ repetition nor
Fourier-null.  A nonbacktracking grammar can delete immediate reversal, but
then the grammar and determinant have changed; it does not rescue the
self-adjoint block.

\begin{theorem}[No-visibility trilemma]
\label{thm:trilemma}
Within SD-C14, a scalar coupling has exactly three possible outcomes:
\begin{enumerate}
 \item it changes some positive trace coefficient and corrupts the Euler
       ledger;
 \item all new trace coefficients cancel, in which case the scalar
       trace-log determinant remains blind to the coupling;
 \item it is detected only by a normalized spectral, magnitude, chiral, or
       other observable that is not the same holomorphic Euler determinant.
\end{enumerate}
No outcome supplies determinant-visible Haar motion while retaining the
frozen all-order ledger.
\end{theorem}

\begin{proof}
By \cref{thm:blindness}, the scalar germ is determined coefficientwise by
all aggregated positive traces.  If a coefficient changes, the ledger
changes.  If none changes, the germ is identical.  The remaining case changes
the observable, as the explicit formulas above demonstrate.
\end{proof}

The frozen adversarial controls sharpen this theorem.
\begin{center}\small
\begin{tabularx}{0.98\textwidth}{L{0.22\textwidth}X L{0.24\textwidth}}
\toprule
Control & Exact outcome & Interpretation\\\midrule
$c=0$ & $W$ reduces to the rigid visible point mass & Paper11 baseline\\
cyclic $N$-fiber & first defect at $r=N$ & finite dimension only delays leakage\\
perturbed Haar density & first nonzero Fourier mode changes that repetition & invisibility is nongeneric\\
normalized trace & moment becomes $1/(1+c)$ & state normalization fails ledger\\
composite/random clocks & same factor $1-za_n$ for every positive inventory & no arithmetic selectivity\\
inverse recurrent edges & balanced word contributes at power two & coupling pollutes grammar\\
\bottomrule
\end{tabularx}
\end{center}

The finite audit realizes these controls without fitted parameters.  For
every $N=2,\ldots,64$, checked through $r=128$, the cyclic control first leaks
at exactly $r=N$.  Frequencies $1,\ldots,16$ at four signed amplitudes leak at
the preregistered frequency in all $64$ cases.  For
$c\in\{0.25,1,3\}$ and twelve complex $q$-values, the largest numerical
residual in the Fuglede--Kadison formula is
$4.440892098500626\times10^{-16}$.  Finally, tensor-prime, composite-only,
and seeded random-increasing inventories, each truncated to $128$ atoms,
give analytic-determinant difference and phase range exactly zero for all
nine inventory--$c$ combinations.  Thus the strongest numerical result is a
\emph{PROVES\_TOO\_MUCH} control, not evidence for a target divisor.

The inventory control is decisive.  For any positive atom weights $(a_n)$ in
the convergence domain,
\[
 \prod_nD_c(za_n)=\prod_n(1-za_n).                          \tag{5.4}
\]
Nothing in \cref{thm:blindness} distinguishes primes from composites or
random increasing clocks.  The mechanism therefore proves too much as an
arithmetic selector.
